\chapter{方法}

\section{总体研究框架}

本研究围绕"自回归滞后特征的捷径竞争是否为气象物理量被稀疏化剪除的主要机制"这一核心问题，设计了四幕式（Four-Act）递进实验框架，如\cref{fig:system-flowchart}所示。

\begin{figure}[htbp]
\centering
\includegraphics[width=\textwidth]{figures/system_flowchart.pdf}
\caption{四幕式递进研究框架系统流程图}
\label{fig:system-flowchart}
\end{figure}

\textbf{第一幕：基线建立（S1）。}以标准KAN训练流程处理净负荷差分预测任务，在多预测时域（$h=6, 12, 24$）上获得教师模型性能基线，并观察各特征组的存活与剪除模式。此阶段确认"气象物理量被系统性剪除"这一现象的存在性和普遍性。

\textbf{第二幕：问题发现（S2）。}在S1基线的基础上，通过消融实验系统检查不同太阳辐照信息源（GHI实测值 vs.\ 天文几何计算值）对模型性能和符号提取结果的影响，进一步锁定问题的范围。

\textbf{第三幕：机制检验（S2-Blocking）。}这是本研究的核心实验阶段。通过控制变量干预——阻断自回归滞后特征的输入通路——直接测量气象物理量可辨识性的变化。阻断前后的对比遵循预注册式的判定规则（详见\ref{sec:blocking-protocol}节），为捷径竞争假说提供干预性证据。

\textbf{第四幕：构造性解决方案（S3）。}基于机制检验的发现，提出结构化分解策略：将净负荷预测分解为负荷、风电、光伏三个物理上自然的子任务，每个子任务仅包含与自身物理机制相关的特征和滞后项，从而在结构层面消解捷径竞争，使气象物理量在各子模型的符号公式中得以保留。

贯穿四幕实验的技术流水线为：数据预处理 $\to$ KAN教师模型训练 $\to$ L1+熵稀疏化 $\to$ 结构化剪枝 $\to$ 符号函数提取 $\to$ SymPy公式构建 $\to$ 预测性能评估与可辨识性诊断。需要强调的是，本研究不修改KAN的训练算法或网络结构——所有干预手段仅体现在输入特征的组成控制（滞后阻断）和任务结构的设计（S3分解）上。这一设计确保了实验结论的归因清晰性：观察到的可辨识性变化只能归因于特征空间或任务结构的改变，而非模型架构的调整。这里所称teacher亦指``KAN teacher + symbolic extraction''的后验提取链路，而不是经典教师-学生蒸馏。

\section{KAN教师模型训练}

\subsection{KAN基本原理}

Kolmogorov-Arnold网络（Kolmogorov-Arnold Network, KAN）\cite{kolmogorov1957representation}是一种基于Kolmogorov-Arnold叠加定理\cite{arnold1957functions}的神经网络架构。不同于传统多层感知机（MLP）在节点上使用固定激活函数、在边上使用线性权重的范式，KAN将可学习的一元函数置于网络的边上，而节点仅执行求和操作。

具体而言，对于一个两层KAN，其计算过程可表示为：
\begin{equation}
f(\mathbf{x}) = \sum_{j=1}^{m} \Phi_{j}^{(2)}\left(\sum_{i=1}^{n} \phi_{ij}^{(1)}(x_i)\right)
\end{equation}
其中$\phi_{ij}^{(1)}$和$\Phi_{j}^{(2)}$均为可学习的一元函数。在pykan实现中\cite{liu2025kan}，每条边上的一元函数以B样条（B-spline）参数化：
\begin{equation}
\phi(x) = w_b \cdot \text{silu}(x) + w_s \cdot \text{spline}(x)
\end{equation}
其中$\text{spline}(x)$是$k$阶B样条基函数的线性组合，$w_b$和$w_s$分别为基础权重和样条权重。

KAN架构的关键优势在于其天然的符号回归（symbolic regression）亲和性：每条边上的一元函数可以独立地被近似为已知的解析函数（如$x^2$、$\sin(x)$、$\exp(x)$等），进而通过符号组合（symbolic composition）将整个网络转化为闭式数学表达式。这一特性使KAN成为可解释预测模型的理想载体。

\subsection{网络架构设计}

本研究采用单隐藏层KAN架构，网络结构表示为$[\text{input\_dim}, \text{hidden\_width}, 1]$。其中：

\begin{itemize}
  \item input\_dim：由特征组配置动态决定，取决于所选特征组的变量数量
  \item hidden\_width：默认为10，即隐藏层包含10个节点
  \item 输出层为单节点，对应标量回归目标
\end{itemize}

B样条参数设置为网格数$G \in \{3, 5\}$、样条阶数$k=3$（三次B样条）。网格范围设为$[-5, 5]$，覆盖Z-score标准化后的典型输入值域。

该架构设计与MCKAN\cite{chen2025mckan}中的思想存在对应关系：本研究中的多尺度滞后特征（lag\_12, lag\_24, lag\_48）在概念上类似于MCKAN的多尺度卷积核，均旨在捕捉不同时间尺度上的序列动态；而特征组的显式划分则与MCKAN中基于Pearson相关系数的特征选择策略相呼应。

\subsection{训练调度}

KAN模型的训练分为四个阶段，每个阶段有明确的优化目标和超参数设置：

\textbf{阶段一：预热训练（Warmup）。}200步，学习率$\eta = 0.01$。目标是让B样条函数获得基本的数据拟合能力，在此阶段不施加稀疏正则化。预热结束后可选择性地更新B样条网格以适应数据分布。

\textbf{阶段二：稀疏化训练（Sparsify）。}800步，学习率$\eta = 0.005$。在损失函数中加入L1范数和熵正则化项：
\begin{equation}
\mathcal{L} = \mathcal{L}_{\text{MSE}} + \lambda \cdot \left(\lambda_{L1} \cdot \|\boldsymbol{\phi}\|_1 + \lambda_{\text{entropy}} \cdot H(\boldsymbol{\phi})\right)
\end{equation}
其中$\lambda = 0.01$为总正则化强度（标准设置），$\lambda_{L1} = 1.0$为L1范数系数，$\lambda_{\text{entropy}} = 2.0$为熵正则系数。L1范数促使多数边的激活函数幅值趋近于零；熵正则鼓励激活幅值的分布趋向稀疏（即少数边保持大幅值，多数边衰减为零），而非均匀分散。两项正则化的联合作用是本研究关注的核心机制：在有限的模型容量下，稀疏化过程形成特征间的竞争，预测能力更强的特征（如自回归滞后项）优先获得激活边的分配，而预测能力相对较弱的特征（如气象物理量）则被压制。

\textbf{阶段三：结构化剪枝（Prune）。}对稀疏化后的模型执行结构化剪枝，移除幅值低于阈值的激活边和非活跃节点。剪枝采用网格搜索策略，设置9组候选阈值（node\_th $\in \{0.01, 0.02, 0.03\}$与edge\_th $\in \{0.002, 0.005, \ldots, 0.20\}$的组合），从中选择满足以下约束的最优方案：
\begin{itemize}
  \item RMSE退化约束：剪枝后RMSE不超过剪枝前的$1.1$倍（即退化率$\leq 10\%$）
  \item 稀疏度目标：边剪除率$\geq 80\%$
\end{itemize}
在满足RMSE约束的候选中，优先选择RMSE最低者；若无候选同时满足两项约束，则采取降级选择策略。

\textbf{阶段四：精细化训练（Refine）。}200步，学习率$\eta = 0.5$。在剪枝后的稀疏网络上继续训练，恢复因剪枝造成的性能损失。较高的学习率有利于在参数空间的受限子流形上快速收敛。

\section{符号提取}

符号提取（symbolic extraction）是将剪枝后KAN教师模型转化为闭式数学表达式的关键步骤，分为边级拟合和全局组合两个阶段。

\subsection{边级符号拟合}

对剪枝后模型中每条存活的激活边$\phi_{ij}^{(l)}$，执行以下符号拟合流程：

\begin{enumerate}
  \item \textbf{候选函数库定义。}本研究使用三组候选函数库：
  \begin{itemize}
    \item strict库：$\{x, x^2, x^3, \sin, \cos, |\cdot|\}$，仅包含基础代数和三角函数
    \item medium库：$\{x, x^2, x^3, \sin, \cos, |\cdot|, \exp\}$，额外包含指数函数
    \item strict\_poly4库：$\{x, x^2, x^3, x^4, \sin, \cos, |\cdot|\}$，在strict库基础上加入$x^4$，用于区分“高阶多项式扩展”和“指数扩展”两类表达能力
  \end{itemize}

  \item \textbf{拟合与评分。}对每个候选函数$g$，pykan内部通过最小二乘法拟合仿射变换参数$a, b, c, d$，使得$a \cdot g(b \cdot x + c) + d$逼近该边的B样条激活值，并计算拟合$R^2$。最终评分综合拟合质量和复杂度：
  \begin{equation}
  \text{score} = R^2 \times w_{\text{simple}}^{\text{complexity}}
  \end{equation}
  其中$w_{\text{simple}} = 0.9$为简洁性偏好权重。

  \item \textbf{固定决策。}若最优候选的$R^2$超过给定阈值，则将该边固定为对应的符号函数。strict与medium使用$\{0.98, 0.99, 0.995\}$三档阈值；strict\_poly4使用更严格的$\{0.99, 0.995, 0.999\}$。

  \item \textbf{零边处理。}对于在剪枝阶段已被移除的边（mask值为零），直接固定为常数函数$0$，确保其不对符号公式产生贡献。
\end{enumerate}

\subsection{全局公式构建}

完成所有边的符号拟合后，通过符号组合（symbolic composition）将各层的一元函数连接为完整的多元表达式。pykan提供的\texttt{symbolic\_formula}方法以特征变量名为输入符号，依次将每层的符号边函数代入下一层的求和结构，最终输出SymPy表达式对象。

在公式构建过程中，输入标准化参数（Z-score的均值和标准差）和输出标准化参数会被嵌入公式，使得最终表达式直接接受原始尺度的输入并输出原始尺度的预测值。

\subsection{安全评估}

符号公式的数值评估引入两项安全机制以防止数值溢出：
\begin{itemize}
  \item \textbf{安全指数截断}（safe\_exp\_clip）：将$\exp(\cdot)$的参数截断至$[-10, 10]$范围，防止指数爆炸
  \item \textbf{分位数截断}（eval\_clip\_quantiles）：在评估时按$(0.01, 0.99)$分位数截断输入特征，消除极端离群值的影响
\end{itemize}

\section{可辨识性诊断指标}

为定量评估KAN符号回归过程中变量的存活与剪除行为，本研究定义了一组可辨识性（identifiability）诊断指标。这些指标构成了阻断实验和结构化分解实验的核心测量工具。

\subsection{变量进入率（Variable Entry Rate, VER）}

变量进入率衡量某一输入变量在多次独立训练（不同随机种子）中被保留在剪枝后网络结构中的频率：
\begin{equation}
\VER_v = \frac{\#\{s \in \mathcal{S} : \text{edge\_count}(v, s) \geq 1\}}{|\mathcal{S}|}
\end{equation}
其中$\mathcal{S}$为随机种子集合，$\text{edge\_count}(v, s)$表示在种子$s$下，变量$v$在剪枝后仍保有的激活边数量。$\VER_v = 1$意味着变量$v$在所有种子下均至少有一条存活边；$\VER_v = 0$则意味着该变量在所有种子下均被完全剪除。

\subsection{公式出现率（Formula Appearance Rate, FAR）}

公式出现率进一步考察变量是否成功通过符号拟合阶段、最终进入SymPy闭式公式：
\begin{equation}
\FAR_v = \frac{\#\{s \in \mathcal{S} : v \in \text{symbols}(f_{\text{sym}}^{(s)})\}}{|\mathcal{S}|}
\end{equation}
其中$\text{symbols}(f_{\text{sym}}^{(s)})$表示种子$s$下最终符号公式中出现的自由变量集合。FAR $\leq$ VER恒成立，因为一个变量在剪枝后存活（VER条件）并不保证其对应的边能被符号函数库中的候选函数以足够高的$R^2$拟合（FAR条件）。

\subsection{激活边计数（Edge Count）}

激活边计数作为VER的连续值补充，记录变量$v$在剪枝后保有的激活边数量：
\begin{equation}
\text{edge\_count}(v) = \sum_{j=1}^{m} \mathbb{1}\left[\text{mask}_{v,j} \neq 0\right]
\end{equation}
其中求和遍历隐藏层的所有节点。

\subsection{迁移间隙比（Transfer Gap Ratio, TGR）}

迁移间隙比度量符号公式相对于教师KAN模型的性能退化：
\begin{equation}
\TGR = \frac{\RMSE_{\text{symbolic}}}{\RMSE_{\text{teacher}}}
\end{equation}
$\TGR = 1$表示符号公式完美复现了教师模型的预测精度；$\TGR > 1$表示符号转换引入了性能损失。

\subsection{差异VER（$\Delta$VER）}

差异VER是阻断实验的核心测量指标，定义为同一变量在阻断条件和非阻断条件下VER的差值：
\begin{equation}
\Delta\VER_v = \VER_v^{(\text{blocked})} - \VER_v^{(\text{unblocked})}
\end{equation}
若$\Delta\VER_v > 0$，说明阻断自回归滞后特征后变量$v$的存活率上升，支持捷径竞争假说。

各诊断指标的定义汇总于\cref{tab:metrics}。

\begin{table}[htbp]
\centering
\caption{可辨识性诊断指标定义}
\label{tab:metrics}
\begin{tabular}{ccccc}
\toprule
指标 & 符号 & 定义 & 取值范围 & 用途 \\
\midrule
变量进入率 & $\VER_v$ & 变量在剪枝后至少保留1条边的种子比例 & $[0, 1]$ & 衡量变量存活稳健性 \\
公式出现率 & $\FAR_v$ & 变量出现在最终公式中的种子比例 & $[0, 1]$ & 衡量符号化成功率 \\
激活边计数 & edge\_count$(v)$ & 变量在剪枝后的存活边数 & $\mathbb{N}_0$ & VER的连续补充 \\
迁移间隙比 & TGR & $\RMSE_{\text{sym}} / \RMSE_{\text{teacher}}$ & $[1, +\infty)$ & 符号提取质量诊断 \\
差异VER & $\Delta\VER_v$ & $\VER^{\text{blocked}} - \VER^{\text{unblocked}}$ & $[-1, 1]$ & 阻断效应核心测量 \\
\bottomrule
\end{tabular}
\end{table}

\section{自回归阻断协议（S2-Blocking）}
\label{sec:blocking-protocol}

自回归阻断实验是本研究的核心机制检验手段。其设计原则为受控干预（controlled intervention）：在保持所有其他条件完全一致的前提下，仅改变输入特征中自回归滞后项的有无，观察气象物理量可辨识性的变化。

\subsection{实验条件设计}

阻断实验设置两组核心对比条件：

\textbf{Case 3：聚焦风电任务。}预测目标为$\Delta\text{wind}_{h6}$，特征组包含cyclic和meteo\_wind。
\begin{itemize}
  \item \textbf{非阻断条件（unblocked）}：lag\_series = [``wind'']，即包含wind\_lag\_24和wind\_lag\_48
  \item \textbf{阻断条件（blocked）}：lag\_series = []，即移除所有滞后特征
\end{itemize}

\textbf{Case 4：直接净负荷任务。}预测目标为$\Delta\text{net\_load}_{h6}$，特征组包含meteorology、solar和cyclic的完整集合。
\begin{itemize}
  \item \textbf{非阻断条件（unblocked）}：lag\_series = [``load'', ``wind'', ``solar'']
  \item \textbf{阻断条件（blocked）}：lag\_series = [``load'']，即仅保留负荷滞后项
\end{itemize}

\subsection{控制条件}

为确保两组条件之间的可比性，以下参数在阻断与非阻断条件之间完全固定：
\begin{itemize}
  \item KAN架构：$[\text{input\_dim}, 10, 1]$，$G=5$，$k=3$
  \item 正则化参数：$\lambda = 0.01$，$\lambda_{L1} = 1.0$，$\lambda_{\text{entropy}} = 2.0$
  \item 训练调度：warmup 200步 $\to$ sparsify 800步 $\to$ prune $\to$ refine 200步
  \item 剪枝策略：相同的9组候选阈值和选择规则
  \item 符号函数库：strict库（$x, x^2, x^3, \sin, \cos, |\cdot|$）
  \item 符号拟合$R^2$阈值：0.99
  \item 随机种子：1, 2, 3, 4, 5（共5个种子）
\end{itemize}

\subsection{预注册判定规则}

为避免事后合理化（post-hoc rationalization），本研究预先指定了阻断实验的判定标准：

\textbf{成功（SUCCESS）}：对至少一个气象物理变量$v$，$\Delta\VER_v$的Bootstrap 95\%置信区间下界$> 0$。Bootstrap采用配对抽样（across same seeds），重采样次数为10,000次，置信区间取百分位法。

\textbf{强成功（STRONG SUCCESS）}：至少2个气象物理变量满足$\Delta\VER_v > 0$的SUCCESS条件。

\textbf{零结果（NULL）}：所有气象物理变量的$\Delta\VER_v$置信区间均包含0。

\subsection{逻辑链条}

阻断实验的推理逻辑可归纳为：
\begin{enumerate}
  \item 若$\Delta\VER_v > 0$（阻断后可辨识性上升），则说明变量$v$在非阻断条件下的低可辨识性并非因为其对目标缺乏预测贡献，而是因为自回归滞后特征在稀疏化竞争中占据了模型容量
  \item 若$\Delta\VER_v \leq 0$，则不能排除变量$v$确实对目标贡献有限的可能性
\end{enumerate}

为避免把teacher/composite predictor与symbolic formula混写成同一层证据，\cref{tab:protocol-ledger}固定了各实验段落的统计单位与对象层级；完整CSV资产另存于 `doc/paper_assets/paper_delivery_20260306/protocol_ledger_20260417.csv`。

\begin{table}[htbp]
\centering
\caption{实验协议与证据对象 ledger}
\label{tab:protocol-ledger}
\begin{tabularx}{\textwidth}{llllX}
\toprule
证据块 & 目标 & lag设置 & 统计单位 & 证据对象 \\
\midrule
5.1 主精度 & $\Delta\text{net\_load}_{h6}$ & load/wind/solar & run & teacher(abs-test replay) \\
5.2 direct symbolic & $\Delta\text{net\_load}_{h6}$ & load/wind/solar & symbolic config & direct symbolic formula \\
5.4 Case 3 & $\Delta\text{wind}_{h6}$ & wind / none & seed & pruned structure + formula stability \\
5.4 Case 4 unblocked & $\Delta\text{net\_load}_{h6}$ & load/wind/solar & symbolic config & direct symbolic baseline \\
5.4 Case 4 blocked & $\Delta\text{net\_load}_{h6}$ & load only & seed & blocked pruned structure \\
5.5 S3 predictor & $\text{net\_load}_{h6}$ & structured combo & run & composite predictor \\
5.5 S3 formula & $\text{net\_load}_{h6}$ & formula replay & run & replayed total formula \\
\bottomrule
\end{tabularx}
\end{table}

其中Case 4最需要单独说明：其unblocked基线来自9组direct symbolic配置，而blocked结果来自5个seed的活跃边统计，因此后文只将该对照解释为阻断支持证据，而不写成同层级的严格闭合比较。

\section{结构化分解策略（S3）}
\label{sec:s3-design}

\subsection{设计动机}

S2-Blocking实验揭示了捷径竞争的存在，但阻断滞后特征意味着牺牲预测精度以换取可解释性。S3结构化分解策略的目标是：在不牺牲预测精度的前提下，通过任务结构的重新设计来缓解捷径竞争，使气象物理量在各子任务的局部符号公式中得到保留。

其设计动机基于净负荷的物理可分解性。净负荷的定义式$\text{net\_load} = \text{load} - \text{wind} - \text{solar}$表明，净负荷的变化由三个物理上独立的分量驱动：
\begin{itemize}
  \item \textbf{负荷变化}$\Delta\text{load}$：主要由气温和人类活动周期驱动
  \item \textbf{风电变化}$\Delta\text{wind}$：主要由风速驱动
  \item \textbf{光伏变化}$\Delta\text{solar}$：主要由太阳辐照度和太阳几何位置驱动
\end{itemize}

\subsection{子任务定义}

S3实验将净负荷差分预测分解为三个子任务，每个子任务由独立的KAN教师模型（sub-KAN）处理。各子任务的特征配置汇总于\cref{tab:s3-config}。

\begin{table}[htbp]
\centering
\caption{S3结构化分解子任务配置}
\label{tab:s3-config}
\begin{tabular}{ccccc}
\toprule
子任务 & 目标变量 & 特征组 & 滞后项 & 核心物理驱动 \\
\midrule
Load sub-KAN & $\Delta\text{load}_{h6}$ & \makecell{cyclic, meteo\_temp,\\meteo\_degree, meteo\_pressure,\\solar\_flag} & load\_lag\_\{12,24,48\} & 气温、度日数 \\
\addlinespace
Wind sub-KAN & $\Delta\text{wind}_{h6}$ & cyclic, meteo\_wind & wind\_lag\_\{24,48\} & 风速、$v^3$、轮毂风速 \\
\addlinespace
Solar sub-KAN & $\Delta\text{solar}_{h6}$ & \makecell{cyclic, solar\_geom,\\solar\_flag, meteo\_irradiance,\\meteo\_temp} & solar\_lag\_\{24,48\} & GHI、太阳高度角 \\
\bottomrule
\end{tabular}
\end{table}

\subsection{组合规则}

三个子模型的预测结果通过固定加法规则（无学习权重）组合为净负荷差分的复合预测：
\begin{equation}
\Delta\widehat{\text{net\_load}}_{\text{composite}} = \Delta\widehat{\text{load}} - \Delta\widehat{\text{wind}} - \Delta\widehat{\text{solar}}
\end{equation}
该组合规则直接来源于净负荷的定义恒等式，不引入额外的可学习参数，确保了组合预测的物理一致性和可解释性。本文当前已经完成$F_{\text{net}} = F_{\text{load}} - F_{\text{wind}} - F_{\text{solar}}$的显式构造与test集逐点复算，因此S3不再停留在"局部公式 + 结构化组合预测器"的半闭环状态；但总公式相对组合预测器仍存在明显transfer gap，故后续分析仍需严格区分预测器对象与公式对象。

\subsection{可辨识性保留标准}

S3实验以"可辨识性保留"（identifiability preservation）作为成功标准。具体定义为：对于每个子任务，其核心物理驱动变量的VER应$\geq 3/5$（即在5个种子中至少有3个种子保持该变量的存活边）。

与S2-Blocking的纯诊断性设计不同，S3是一个构造性解决方案：它不仅证明了捷径竞争机制的存在，还通过结构化的任务分解提供了一条在保持预测精度的同时恢复物理可解释性的技术路径。

\section{本章小结}

本章详细阐述了研究方法的完整体系。首先介绍了四幕式递进研究框架（基线 $\to$ 问题发现 $\to$ 机制检验 $\to$ 构造性解决方案）的设计逻辑。在技术层面，依次描述了KAN教师模型的训练调度（预热 $\to$ 稀疏化 $\to$ 剪枝 $\to$ 精细化）、符号提取的边级拟合与全局组合流程、以及安全评估机制。

在实验设计层面，定义了五个可辨识性诊断指标（VER、FAR、edge\_count、TGR、$\Delta$VER），为定量评估变量的存活行为提供了标准化工具。核心的S2-Blocking自回归阻断协议通过控制变量干预和预注册判定规则，为捷径竞争假说的干预检验提供了严谨的实验框架。最后，S3结构化分解策略将净负荷预测任务分解为三个物理自然的子任务，旨在从结构层面消解捷径竞争并恢复气象物理量的可辨识性。

所有方法的共同设计原则是：不修改KAN的训练算法或网络架构，仅通过控制输入特征组成和任务结构来实施干预，确保实验结论的归因清晰性。
